% optimization/alignment/60_parametric_and_hybrid.tex

\section{Parametric Representation instead of Sampling}

\subsection{Motivation}

Classical discretization represents curvature and cant by samples:
\[
\kappa_i \approx \kappa(s_i),
\qquad
\phi_i \approx \phi(s_i).
\]

This leads to high-dimensional optimization problems with weak
structure and limited interpretability.

Within the AIM framework, curvature and cant are instead represented
by parameterized functions.


\subsection{Parametric Curvature Model}

\begin{definition}[Parametric curvature]
A transition is described by:
\[
\kappa(s)
=
\kappa_0 + (\kappa_1 - \kappa_0)
\widehat{\kappa}\!\left(\frac{s}{L}, \mathbf{T}\right),
\]
where:
\begin{itemize}
\item \(L\) is the element length,
\item \(\mathbf{T}\) is a vector of transition parameters,
\item \(\widehat{\kappa}\) is a normalized prototype function.
\end{itemize}
\end{definition}


\subsection{Parametric Cant Model}

\begin{definition}[Parametric cant]
Cant is represented analogously:
\[
\phi(s)
=
\phi_0 + (\phi_1 - \phi_0)
\widehat{\phi}\!\left(\frac{s}{L}, \mathbf{T}_\phi\right).
\]
\end{definition}


\subsection{Parameter Vector}

\begin{definition}[Reduced parameter vector]
The optimization variable becomes:
\[
\mathbf{x}
=
(L, \kappa_0, \kappa_1, \mathbf{T},
 \phi_0, \phi_1, \mathbf{T}_\phi).
\]
\end{definition}

The dimension is independent of discretization density.


\subsection{Dimensional Reduction}

A sampled representation uses \(O(N)\) variables.

The parametric representation uses:
\[
O(1) \text{ parameters per element}.
\]

This yields a drastic reduction in problem size.


\subsection{Derivative Structure}

\begin{proposition}
Derivatives reduce to:
\[
\frac{\partial \kappa}{\partial p}
=
(\kappa_1 - \kappa_0)
\frac{\partial \widehat{\kappa}}{\partial p}.
\]
\end{proposition}

Thus, all derivatives are computed analytically from prototype
functions.


\subsection{Integration into Reconstruction}

The parametric representation is evaluated through:
\[
\theta(s) = \int \kappa(s)\,\dd s,
\qquad
\gamma(s) = \int (\cos\theta,\sin\theta)\,\dd s.
\]

This yields a smooth and consistent geometry without discretization
artifacts.


\subsection{Comparison to Sample-Based Methods}

\begin{center}
\begin{tabular}{lcc}
\toprule
 & Sample-based & Parametric \\
\midrule
dimension & high & low \\
smoothness & enforced numerically & intrinsic \\
interpretability & low & high \\
constraints & many local & few global \\
\bottomrule
\end{tabular}
\end{center}


\subsection{Interpretation}

The parametric approach shifts the problem from:
\begin{itemize}
\item fitting values
\end{itemize}
to:
\begin{itemize}
\item designing functions.
\end{itemize}


\subsection{Connection to Transition Families}

Different transition curves correspond to different prototype functions
\(\widehat{\kappa}\).

Thus, the optimization problem includes selection and tuning of
transition families.


\subsection{Computational Implications}

\begin{itemize}
\item smaller optimization problems,
\item better conditioning,
\item faster convergence,
\item direct integration with analytical Jacobians.
\end{itemize}


\subsection{Connection to Classical Systems}

This formulation is conceptually close to classical systems such as
AXTRAN, which operate on element parameters rather than sampled data.


\subsection{Connection to AIM}

\begin{summary}
The parametric representation replaces discrete sampling by structured
functions.

It enables a compact, interpretable, and analytically differentiable
model, which forms the basis for efficient and robust alignment
optimization.

This approach is central to the AIM philosophy.
\end{summary}


\section{Hybrid Parametric--Discrete Model}

\subsection{Motivation}

Purely parametric models provide low-dimensional structure but may lack
flexibility for fitting real-world data.

Purely sampled models provide flexibility but lead to large and
ill-conditioned optimization problems.

A hybrid approach combines both advantages.


\subsection{Hybrid Representation}

\begin{definition}[Hybrid curvature model]
\[
\kappa(s)
=
\kappa_{\mathrm{param}}(s;\mathbf{T})
+
\delta\kappa(s),
\]
where:
\begin{itemize}
\item \(\kappa_{\mathrm{param}}\) is a parametric base model,
\item \(\delta\kappa(s)\) is a correction term.
\end{itemize}
\end{definition}

\begin{definition}[Hybrid cant model]
\[
\phi(s)
=
\phi_{\mathrm{param}}(s;\mathbf{T}_\phi)
+
\delta\phi(s).
\]
\end{definition}


\subsection{Discretization of Correction}

The correction term is represented discretely:
\[
\delta\kappa_i = \delta\kappa(s_i),
\qquad
\delta\phi_i = \delta\phi(s_i).
\]


\subsection{Extended Parameter Vector}

\begin{definition}
\[
\mathbf{x}
=
(\mathbf{T}, \mathbf{T}_\phi, \delta\boldsymbol{\kappa}, \delta\boldsymbol{\phi}).
\]
\end{definition}


\subsection{Regularization of Corrections}

To avoid overfitting, the correction terms are penalized:
\[
J_{\delta}
=
\int_0^L (\delta\kappa'(s))^2 \,\dd s
+
\int_0^L (\delta\phi'(s))^2 \,\dd s.
\]

This enforces smooth, small deviations from the parametric model.


\subsection{Residual Structure}

The residual now consists of:
\begin{itemize}
\item parametric contributions,
\item correction contributions,
\item data fitting terms.
\end{itemize}


\subsection{Block Structure}

The parameter vector decomposes as:
\[
(\mathbf{T}, \delta\boldsymbol{\kappa})
\quad \text{and} \quad
(\mathbf{T}_\phi, \delta\boldsymbol{\phi}).
\]

This yields a multi-level block structure:
\begin{itemize}
\item coarse scale (parameters),
\item fine scale (corrections).
\end{itemize}


\subsection{Interpretation}

The parametric model captures:
\begin{itemize}
\item geometric intent,
\item transition structure,
\item engineering semantics.
\end{itemize}

The correction term captures:
\begin{itemize}
\item measurement noise,
\item local deviations,
\item modelling imperfections.
\end{itemize}


\subsection{Connection to Measurement Data}

Measured polylines can be fitted by minimizing:
\[
\sum_i \|\gamma(s_i) - p_i\|^2,
\]
with the hybrid representation.

The parametric part explains the global structure, while the correction
absorbs residual errors.


\subsection{Computational Strategy}

\begin{itemize}
\item first optimize parametric variables,
\item then refine using correction terms,
\item optionally alternate between both levels.
\end{itemize}


\subsection{Relation to LOD}

The hybrid model naturally supports level-of-detail concepts:
\begin{itemize}
\item coarse LOD: parametric model only,
\item fine LOD: parametric + corrections.
\end{itemize}


\subsection{Connection to AIM}

\begin{summary}
The hybrid representation combines structured modelling and local
flexibility.

It enables robust alignment reconstruction from imperfect data while
preserving the interpretability of parametric models.

This approach bridges engineering design and data-driven fitting within
the AIM framework.
\end{summary}
